\section{Technical Proofs and Formalization}
\label{app:proofs}

\paragraph{Economic distance and task recovery.}
For loadings $b_i,b_j$ on centered common shocks $q$ with finite second
moments, the root mean square of $(b_i-b_j)^\top q$ measures the difference
in return responses in units of return. Writing $\Omega$ for the shock
covariance and $a_i=\Omega^{1/2}b_i$ gives the distance $\norm{a_i-a_j}$.
\lean{FreeFoundations.Score.response_sq}
Perfect marginal forecasts can be consistent with different joint responses:
two fair binary returns can agree in one regime and move oppositely in
another while each retains the same marginal distribution.
\lean{FreeFoundations.Alignment.marginals_free_of_state}
\lean{FreeFoundations.Alignment.joint_depends_on_state}

Suppose ideal task predictions satisfy $\tau=Ca+v$, where $a$ is a
risk-scaled exposure and $v$ the mismatch. For a finite-dimensional
exposure space, a linear map satisfying $\norm{Cu}\geq\kappa\norm u$
in all directions with $\kappa>0$ admits a linear decoder $L$ with
\[
  \norm{L\widehat\tau-a}\leq
  (\norm{\widehat\tau-\tau}+\norm v)/\kappa .
\]
Thus the bound tightens as the tasks distinguish exposure directions more
sharply. For finite labels and contexts with positive predicted
probabilities, excess log loss measures average divergence from the task's
conditional label law.
\lean{FreeFoundations.Recovery.exists_decoder}
\lean{FreeFoundations.Recovery.exposure_error_mismatch}
\lean{FreeFoundations.Score.pretraining_excess}

\paragraph{Recoverability and the transformation class.}
If $s=r\circ z$, histories in the same fibre of $z$ share the same
target. Conversely, define $r$ on $z(H)$ by the unique target of each
fibre and use any value in $S$ outside the image. In the chord example
the four vectors have coordinates $(x,y)$ with $y=1,0,4,5$ respectively.
The stretch takes $x=1$ to $x=4$. For heights $y,z$ the squared cosine is
\[
  c^2=\frac{(x^2+yz)^2}{(x^2+y^2)(x^2+z^2)}.
\]
All dot products are positive, so the unit-vector chord is
$\sqrt{2-2\sqrt{c^2}}$. For $AB$ the squared cosine rises from $1/2$
to $16/17$, while for $AC$ it stays at $25/34$. This reverses the
first firm's choice even on the three-firm roster. With the fourth
firm present, $C$ and $D$ select one another before and after the stretch.
\lean{FreeFoundations.Neighbours.chord_neighbour_flip}
Composing the raw-state readout with the inverse stretch preserves
predictions. For a
restricted readout class $\mathcal R$, a bijection $g$ preserves
the class when $\{r\circ g^{-1}:r\in\mathcal R\}=\mathcal R$. A
restricted class preserved under both $g$ and $g^{-1}$ has the same
attainable functions after the recoding. Equality of the transformed
class gives closure under composition and inversion. The fibre
argument is pointwise. A measurable or approximate recovery claim
requires additional assumptions on the spaces and readout class.

Although recoding can change peers, an estimated exposure can
preserve the selected peers when the gap at the neighbourhood
boundary is sufficiently wide. Suppose $a_i$ is a true exposure,
$\widehat a_i$ its estimate, and
$\norm{\widehat a_j-a_j}\leq\delta_j\leq\delta_{\max}$ for every
issuer. Each distance in row $i$ can move by at most
$\delta_i+\delta_{\max}$. Therefore an estimated strict gap between
selected and unselected neighbours exceeding
$2(\delta_i+\delta_{\max})$ preserves that set. For the study's
ten-neighbour rule the relevant gap is between the tenth
and eleventh distances. The certificate applies to estimated exposure
distances when the exposure-error bounds are available.

\paragraph{Statistical alignment and transfer.}
Under squared loss the economic conditional mean must be a measurable
function of the task conditional mean almost surely for every representation
attaining the latter to attain the former.
\lean{FreeFoundations.Sufficiency.sufficient_iff}
Transporting a source-law task error to issuers requires further premises.
If the issuer law is bounded by $\omega$ times the source law, source
root-mean-square error enters the recovery bound multiplied by
$\sqrt{\omega}$. Changes in conditional task means and task--exposure
mismatch enter separately, all divided by the positive separation margin.
\lean{FreeFoundations.Transfer.source_target_recovery}

\paragraph{Multiplier and covariance.}
\begin{proposition}[Spatial stability]
  \label{prop:multiplier}
  Let $K=K^\top\geq 0$ entrywise, with positive degrees $d_i=\sum_j
  K_{ij}$, and set $W=D^{-1}K$ for $D=\operatorname{diag}(d_i)$. For
  $|\rho|<1$, the inverse $P=(I-\rho W)^{-1}$ exists and
  \[
    \norm{P}_2\leq \frac{\sqrt{d_{\max}/d_{\min}}}{1-|\rho|}.
  \]
  On the same roster, let $\widehat W$ have multiplier $\widehat P$
  at the same $\rho$. If both inverse norms are at most $M$ and the
  common innovation covariance has $\norm V_2\leq v$, then
  \[
    \norm{\widehat P V\widehat P^\top-PVP^\top}_2
    \leq 2v|\rho|M^3\norm{\widehat W-W}_2 .
  \]
\end{proposition}
Symmetry makes $d$ stationary for $W$. Jensen's inequality gives
$\norm{Wx}_D\leq\norm{x}_D$, where $\norm{x}_D^2=x^\top Dx$. Hence
$\norm{(I-\rho W)x}_D\geq(1-|\rho|)\norm{x}_D$.
In finite dimension this implies invertibility, and comparing the
weighted and Euclidean norms gives
Proposition~\ref{prop:multiplier}'s first inequality. A symmetric
affinity with isolates gives a block diagonal operator with a zero
block on the isolates, whose multiplier is identity. The bound
applies on the positive-degree block.

The inverse identity gives $\widehat P-P=\rho \widehat P(\widehat
W-W)P$. Expanding the covariance difference as
$(\widehat P-P)V\widehat P^\top+PV(\widehat P-P)^\top$
and using submultiplicativity gives the second inequality. A fixed
portfolio with weights $w$ multiplies the covariance error bound by
$\norm{w}_2^2$. The displayed bound holds with common $\rho$ and $V$.
Changes to those quantities contribute additional terms.

\paragraph{Formalized components.}
The fibre criterion and coding invariance have Lean formalizations.
\lean{FreeFoundations.Alignment.state_optimal_iff}
\lean{FreeFoundations.Geometry.sufficient_equiv}
Lean also checks the neighbour certificate, multiplier bound, and
covariance bound.
\lean{FreeFoundations.Neighbours.neighbourSet_of_estimated_margin}
\lean{FreeFoundations.ReversibleNetwork.normalize_resolvent_l2_bound}
\lean{FreeFoundations.Stability.matrix_covariance_sub_le}
The restricted class group argument and isolate decomposition are
elementary written arguments.

\section{Data Construction and Estimation Details}
\label{app:data}

We weight selected classes by their prior calendar-month
capitalization. CRSP CIZ monthly total returns already incorporate
reported delisting returns. We compound the three observed months of
a quarter once, keeping incomplete quarters unavailable. When a
return row is missing, the absence of a separate terminal-event
flag leaves the cause of exit and its proceeds unknown. The return estimates
therefore describe the observed quarters of the formation cohort.
Formation membership is fixed at the December ranking cutoff. An
issuer retains its formation identity when a selected security's
subsequent source issuer changes. Missing returns and missing
prior-month capitalization are explicit. SEC information uses filing
timestamps before the preceding CRSP close. Negative book equity
where the ratio calls for a positive denominator, stale balance
sheets, and undefined paired history correlations each have separate
availability rules.

The characteristic reimplementation estimates OLS, spatial lag by
instrumental variables and maximum likelihood, and spatial error by
maximum likelihood. Fit states distinguish available estimates from
invalid or boundary fits. For the primary formation, the stationary
quarter bootstrap resamples common date blocks, recomputes betas and
medians, and refits models while the formation graph stays fixed
before filtering unavailable issuers and renormalizing its rows. A
bootstrap distribution for a constrained refit can still be computed
when the corresponding original fit is unavailable. In that case the
percentiles describe the constrained refits, while the
original-sample estimate retains its recorded fit status.

\paragraph{Frozen model readouts.}
For TabPFN 3.5 we take the native query embedding of a single regression
ensemble member. For TabFM 1.1 we extract query rows after the in-context
predictor's final layer normalization \citep{kong2026tabfm}. For
Kumo-Tabular we take the input to the in-context prediction head, and
for LimiX-2 the input to the regression decoder
\citep{nvidia2026kumo,zhang2026limix}. When financial inputs span several
column blocks, we concatenate the issuer's vectors in block order before
computing chord distances.

\paragraph{Coverage of characteristic histories.}
A neutral distance enters a characteristic comparison when fewer than two
annual observations overlap or the paired ranks have undefined correlation.
Coverage varies across characteristics and over the panel
(Table~\ref{tab:neutral}). For ten-neighbour selection we break distance
ties in the ordered issuer roster before joining the two directed choices.

\begin{table}[t]
  \centering
  \fftabletype
  \IfFileExists{generated/tables/matched-neutral.tex}{\input{generated/tables/matched-neutral}}{\ffmissing}
  \caption{Share of issuer pairs assigned neutral distance (\%) in each
    characteristic. The pooled share counts unordered pairs within quarter
  across the matched panel.}
  \label{tab:neutral}
\end{table}

\section{Scope of the Empirical Evidence}
\label{app:scope}

Across issuers in the fixed cohort, we relate median excess returns to a
characteristic-weighted average of peer returns. On the repeated-formation
panel, we compare the likelihood of the observed quarterly returns under
different peer weights. The January 2016 geometry analysis then returns to
the formation roster to examine which financial characteristics organize
firms within the ratio-based embedding.

\begin{table*}[t]
  \centering
  \fftabletype
  \IfFileExists{generated/tables/quarterly-likelihood.tex}{%
    \input{generated/tables/quarterly-likelihood}%
  }{\ffmissing}
  \caption{In-sample lag log-likelihood differences per quarter.
    All graphs use ten-neighbour weights. Ratios denote Fernández
    histories and XBRL denotes CompanyFacts inputs. Positive values
    favour the network in the first column.
    Percentiles resample paired quarterly contributions conditional on
  fitted parameters. The last column adjusts for multiple comparisons.}
  \label{tab:quarterly-comparisons}
\end{table*}

\begin{table}[t]
  \centering
  \fftabletype
  \begin{tabular}{lcc}
    \toprule
    Intercept & $\widehat\rho$ & Refit percentiles \\
    \midrule
    Excluded & \ffnum{quarterly-fernandez-joint-rho} &
    $[\ffnum{quarterly-fernandez-joint-rho-bootstrap-low},\
    \ffnum{quarterly-fernandez-joint-rho-bootstrap-high}]$
    \\
    Included & Invalid or boundary & --- \\
    \bottomrule
  \end{tabular}
  \caption{Joint characteristic kernel lag fit for the December 2015
    US formation.
    Percentiles are the central 95\% of
    \ffnum[0]{quarterly-fernandez-joint-rho-bootstrap-draws}
    stationary-quarter refits. They describe refit variation
    conditional on the formation graph and the observed-quarter design.
  Each draw uses the issuers with sufficient observed returns for that refit.}
  \label{tab:us-benchmark}
\end{table}
\paragraph{Does the encoder add to its inputs?}
The earlier quarterly comparison in Table~\ref{tab:quarterly-comparisons}
reports differences per quarter. Tables~\ref{tab:matched-factors}
and~\ref{tab:matched-encoders} report differences per issuer-quarter on
the common factor sample. With symmetrized ten-neighbour weights, the
ratio-based TabPFN graph gains
4.98 log-likelihood units per quarter over the joint characteristic graph,
while selecting peers directly from the same financial inputs gains 5.91.
The paired encoder-minus-direct difference is therefore $-0.94$ per quarter,
with a 95\% interval of $[-2.67,0.75]$. The positive encoder contrast
against the joint graph has a Holm-adjusted $p$-value of 0.135
(Table~\ref{tab:quarterly-comparisons}).

\paragraph{The characteristic benchmark.}
The joint characteristic kernel has
\ffnum[1]{quarterly-fernandez-joint-effective-neighbours} effective
neighbours per row on average, compared with $n-1$ for a fitted roster of
$n$ issuers under uniform
off-diagonal weights. When most issuers receive similar weights,
firms share much of the same peer return. These averages carry little
cross-sectional variation relative to a more concentrated graph. On
the \ffnum[0]{quarterly-primary-issuers} fitted issuers, the
joint-kernel lag maximum-likelihood estimate without an intercept is
$\widehat{\rho}=\ffnum{quarterly-fernandez-joint-rho}$.
Table~\ref{tab:us-benchmark} records the contrast and its
uncertainty. With an intercept, the joint-kernel fit is flagged as
invalid or on the boundary. The no-intercept estimate is consistent
with the spatial term absorbing much of the common mean in this US sample.

The point fit and a separate profile-likelihood calculation agree on
the saved original-sample betas, issuer medians and reconstructed
graph. The percentile endpoints recompute betas and medians in common
time blocks and restrict the graph when firms drop out of a refit.
They quantify a different source of variation from the standard error
conditional on the generated beta. These diagnostics
characterize the US reimplementation, where the original sample is unavailable.

\paragraph{Historical ratio-only geometry.}
To identify which financial histories organize the learned peers, we
form principal coordinates from chord distances of the January
2016 TabPFN states using only the four ratio histories. Because the
embedding uses information available at formation, the analysis
includes all \ffnum[0]{quarterly-primary-roster} issuers on the
analytical roster, including firms with incomplete subsequent
returns. For each ratio we use the latest
available annual slot and restrict the comparison to its observed
issuers. The leading coordinate has stronger absolute rank correlation with
market capitalization over assets and market-to-book than with
dividend yield or debt maturity. The same ordering holds for explained
rank variation after regression on the leading three coordinates and
an intercept. This subspace comparison is invariant to sign changes and
orthogonal rotations.

The valuation ratios are part of the encoder inputs. Their alignment
describes how the model organizes its own inputs on this date.
The historical ratio embedding puts about 45\% of positive inertia
in its leading coordinate. For the market-beta comparison we restrict
these same coordinates to issuers with an available return-based beta
estimate and regress beta ranks on the three coordinates and an intercept.
On this formation date the learned geometry is closely associated
with a valuation measure supplied to the encoder, while its alignment
with independently estimated market exposure is weak.

\section{Selected broad-input geometry}
\label{app:ablation}

For the broad-input model selected by its market-controlled gain, we rebuild
the peer graph using the first one or three principal coordinates at each
formation. We also omit each engineered ratio history in turn from both
the historical examples and the issuers being encoded. The underlying broad
filing fields remain available, so the omission measures the incremental
role of the supplied ratio history. Each comparison uses the same return
observations and the full broad-input graph as its reference.

\begin{table*}[t]
  \centering
  \fftabletype
  \IfFileExists{generated/tables/matched-ablation.tex}{\input{generated/tables/matched-ablation}}{\ffmissing}
  \caption{Broad-input TabPFN peer graphs after projection or omission.
    Overlap averages issuer-level Jaccard similarities within each quarter
    and then across quarters. Fit differences are nats per issuer-quarter
  relative to the full graph, with paired conditional 95\% intervals.}
  \label{tab:ablation}
\end{table*}

Using three coordinates improves on the first-coordinate projection, although
both projected graphs fit less well than the full graph
(Table~\ref{tab:ablation}). Among the four omissions, removing the
market-equity-to-assets history
changes the peer sets most, even with the broad filing fields and historical
context labels still available. The estimated changes in fit are small,
with all four omission intervals spanning zero.

\section{Matched factor specifications}

\paragraph{Broad-input encoder versus direct features.}
Table~\ref{tab:matched-broad-factors} compares TabPFN with the direct graph
built from the same broad financial disclosures. Both graphs use the common
issuer-quarter mask and symmetrized ten-neighbour weights. The historical
contexts and extracted representations remain fixed. Each factor
specification refits the spatial-lag model; paired stationary-quarter
intervals condition on its fitted loadings and parameters. These additional
factor comparisons do not change the market-controlled encoder selection.

\begin{table}[t]
  \centering
  \fftabletype
  \IfFileExists{generated/tables/matched-broad-factors.tex}{\input{generated/tables/matched-broad-factors}}{\ffmissing}
  \caption{Broad-input TabPFN minus the direct-feature graph in spatial-lag
    log likelihood, in nats per issuer-quarter. Positive gains favour
    TabPFN. All three controls use identical issuer-quarters and peer
    graphs, with paired conditional 95\% intervals. Dashes denote
  unavailable fits.}
  \label{tab:matched-broad-factors}
\end{table}

The following tables report the factor-only, spatial-lag and spatial-error
fits on the same issuer-quarters. Reading down each table shows how fit
changes as factors enter separately, alongside the market, or jointly.
Within a spatial specification, comparing the two peer graphs at the same
factor controls isolates the choice of distance. The parameter counts and
information criteria refer to the return fit conditional on the estimated
issuer exposures.
\newcommand{\ffFactorDetail}[2]{
  \begin{table*}[p]
    \centering
    \fftabletype
    \input{generated/tables/matched-detail-#1}
    \caption{#2 on the common issuer-quarter mask. Information criteria
    condition on estimated first-stage exposures.}
\end{table*}}
\IfFileExists{generated/tables/matched-detail-none.tex}{
  \ffFactorDetail{none}{Factor-only models}
  \ffFactorDetail{lag}{Spatial lag models}
  \ffFactorDetail{error}{Spatial error models}
}{\ffmissing}
